% The LaTeX kernel, ported from latex.ltx rather than invented.
%
% `prelude.tex` is the stand-in: a macro that would have drawn something is
% answered with its text, and a macro whose effect the stomach would have read
% is answered with nothing. This file is the opposite kind of thing. Every
% definition here is the one `latex.ltx` makes, carried over as TeX, with the
% place it came from written above it. Where a definition could NOT be carried
% over, the substitution is named and the primitive that is missing is named
% with it, because "this is not what LaTeX does" is the fact a reader of this
% file most needs.
%
% It is read AFTER `prelude.tex`, so a name that appears in both ends up with
% the ported definition. The stand-in stays for the names this file does not
% reach yet.
%
% The line numbers are against
% /usr/local/texlive/2026/texmf-dist/tex/latex/base/latex.ltx, LaTeX2e
% <2025-11-01>. They move between releases; the definitions do not.
\catcode`\{=1 \catcode`\}=2 \catcode`\#=6 \catcode`\@=11
%
% ---- ltdefns: the plumbing every other definition rests on ---------------
% latex.ltx:1225-1226. A name built out of characters, defined and used.
\def\@namedef#1{\expandafter\def\csname #1\endcsname}
\def\@nameuse#1{\csname #1\endcsname}
% latex.ltx:1090-1091, 896-897.
\def\@empty{}
\long\def\@gobbletwo#1#2{}
\long\def\@gobblefour#1#2#3#4{}
\long\def\@firstofone#1{#1}
\long\def\@firstofthree#1#2#3{#1}
\long\def\@secondofthree#1#2#3{#2}
\long\def\@thirdofthree#1#2#3{#3}
% latex.ltx:1382-1383. \@car takes the head of a list, \@cdr the tail.
\def\@car#1#2\@nil{#1}
\def\@cdr#1#2\@nil{#2}
\def\@nnil{\@nil}
% latex.ltx:1389. Both arguments expanded once, then the command applied.
\def\@expandtwoargs#1#2#3{\edef\reserved@a{\noexpand#1{#2}{#3}}\reserved@a}
% latex.ltx (ltcntrl): the loop the option lists and the counter reset lists are
% both walked with. Carried verbatim -- the delimited parameter texts ARE the
% mechanism, and any paraphrase of them is a different macro.
\long\def\@fornoop#1\@@#2#3{}
\long\def\@for#1:=#2\do#3{%
  \expandafter\def\expandafter\@fortmp\expandafter{#2}%
  \ifx\@fortmp\@empty \else
    \expandafter\@forloop#2,\@nil,\@nil\@@#1{#3}\fi}
\long\def\@forloop#1,#2,#3\@@#4#5{\def#4{#1}\ifx #4\@nnil \else
       #5\def#4{#2}\ifx #4\@nnil \else#5\@iforloop #3\@@#4{#5}\fi\fi}
\long\def\@iforloop#1,#2\@@#3#4{\def#3{#1}\ifx #3\@nnil
       \expandafter\@fornoop \else
      #4\relax\expandafter\@iforloop \fi#2\@@#3{#4}}
%
% ---- counters: ltcounts.dtx ----------------------------------------------
% A LaTeX counter IS a \count register: \@definecounter says
%
%     \expandafter\newcount\csname c@foo\endcsname
%
% and every read of it is `\csname c@foo\endcsname' used as a number. texrs has
% \count, \advance, \divide, \multiply and \the\count, and all of them work --
% but not through a name \csname built, and not through anything that has to be
% FROZEN. Four refusals, each measured against target/debug/texrs:
%
%   \expandafter\countdef\csname c@foo\endcsname=3   assigns, and the matching
%       \expandafter\the\csname c@foo\endcsname      answers `! Unsupported \the.'
%   \count\csname cr@foo\endcsname=0                 works at the top level and
%       answers `! Missing number, found \the.' from inside a macro body: the
%       \csname is not expanded while a register number is being scanned.
%   \edef\x{\csname c@foo\endcsname}                 stores the three tokens
%       rather than the counter's digits: \edef does not expand \csname.
%   \count255=5 \edef\x{\the\count255}\count255=9 \message{\x}   says 9, not 5:
%       \edef does not expand \the, \number or \numexpr either. NOTHING can be
%       frozen into a macro body here, which is what rules out carrying a
%       counter's value as a macro instead of as a register, and what rules out
%       writing plain TeX's \newcount allocator, whose whole job is to freeze
%       the next free register number into a name.
%
% So the value lives in a register, as it does under LaTeX, and what is carried
% by name is the register's NUMBER: `\cr@foo' expands to the digits of the
% register holding counter foo. Reaching the register from a counter name is
% then \expandafter through a DELIMITED parameter --
%
%     \def\@stepctr#1\@nil{\advance\count#1 by 1 }
%     \def\stepcounter#1{\expandafter\@stepctr\csname cr@#1\endcsname\@nil}
%
% -- rather than `\count\csname cr@#1\endcsname', which is the form that fails
% inside a macro. The \@nil delimiter is what makes a two-digit register number
% arrive whole: texrs's \expandafter expands the \csname AND the macro it names,
% so `\cr@foo' arrives as the characters `1' `2' `0' and an undelimited #1 would
% take only the first of them.
%
% The register NUMBERS are allocated in Rust, in `src/latex/counters.rs', and
% reach a run as `\@namedef{cr@NAME}{NN}' lines after this file. A counter with
% no register stops the document at `! Missing number, found \cr@NAME.', which
% is loud and names the counter.
%
% Registers 100..119 are the counters every LaTeX document has; 120 and up are
% the ones a document declares; 254 and 255 are scratch.
%
% latex.ltx: \def\@definecounter#1{...}
\def\@definecounter#1{%
  \setcounter{#1}{0}%
  \expandafter\let\csname cl@#1\endcsname\@empty
  \@addtoreset{#1}{@ckpt}%
  \expandafter\let\csname p@#1\endcsname\@empty
  \@namedef{the#1}{\arabic{#1}}}
% latex.ltx: \def\newcounter#1{...\@ifnextchar[{\@newctr{#1}}{}}
% The optional argument is the counter this one is reset by.
\def\newcounter#1{\@definecounter{#1}\@ifnextchar[{\@newctr{#1}}{}}
\def\@newctr#1[#2]{\@addtoreset{#1}{#2}}
% latex.ltx: \def\@addtoreset#1#2{\expandafter\@cons\csname cl@#2\endcsname{{#1}}}
% \@cons appends `\@elt{#1}' to the list macro; the list is walked with \@elt
% \let to whatever the walk wants done.
\def\@cons#1#2{\begingroup\let\@elt\relax\xdef#1{#1\@elt #2}\endgroup}
\def\@addtoreset#1#2{\expandafter\@cons\csname cl@#2\endcsname{{#1}}}
\let\@elt\relax
\expandafter\let\csname cl@@ckpt\endcsname\@empty
% The three commands that WRITE a counter go through one macro per counter --
% `\ctr@NAME' -- rather than through `\cr@NAME' and the delimited parameter the
% reads use. The reason is what happens to a name NOTHING has allocated.
%
% `\setcounter{a}{1}' for an undeclared `a' expanded to the characters
% `\count \cr@a =1', and an assignment is not expandable, so in every place a
% macro is expanded WITHOUT being executed -- a \message body above all -- those
% characters were written out as text. Measured: `\message{[\setcounter{a}{1}]}'
% said `[\count \cr@a =1 ]' where the stand-in this file replaced said `[]', and
% `tests/latex.rs', state_macros_consume_their_arguments_and_produce_nothing,
% pins the empty answer.
%
% There is no way for a macro to know it is inside a \message, so the branch is
% taken by NAME instead, where Rust can take it: `src/latex/counters.rs' writes
% `\ctr@NAME' for every counter name the source mentions -- the real assignment
% for a name with a register, and a two-argument macro producing NOTHING for a
% name without one, which is LaTeX's `! No counter `a' defined.' as far as an
% engine with no error recovery can carry it. A counter that HAS a register is
% still an assignment and is still written out inside a \message; the counters
% every document actually steps are declared, and the value is right everywhere
% it is executed.
%
% \ctr@NAME takes the operation and the operand: `\ctr@section\@ctrset{4}'.
\def\@ctrset#1#2{\count#1=#2 }
\def\@ctradd#1#2{\advance\count#1 by #2 }
% latex.ltx: \def\setcounter#1#2{...\global\csname c@#1\endcsname#2\relax}
\def\setcounter#1#2{\@nameuse{ctr@#1}\@ctrset{#2}}
% latex.ltx: \def\addtocounter#1#2{...\global\advance\csname c@#1\endcsname#2\relax}
\def\addtocounter#1#2{\@nameuse{ctr@#1}\@ctradd{#2}}
% latex.ltx: \def\stepcounter#1{\addtocounter{#1}\@ne
%              \begingroup \let\@elt\@stpelt \csname cl@#1\endcsname \endgroup}
% The second half is what makes \section reset \subsection: the counters listed
% in \cl@section are each set back to zero.
\def\stepcounter#1{%
  \@nameuse{ctr@#1}\@ctradd{1}%
  \begingroup \let\@elt\@stpelt \csname cl@#1\endcsname \endgroup}
% latex.ltx: \def\@stpelt#1{\global\csname c@#1\endcsname \m@ne\stepcounter{#1}}
% Set to -1 and step, so a counter with subordinates of its own resets those too.
\def\@stpelt#1{\setcounter{#1}{-1}\stepcounter{#1}}
% latex.ltx: \def\refstepcounter#1{\stepcounter{#1}%
%              \protected@edef\@currentlabel{\csname p@#1\endcsname\csname the#1\endcsname}}
% This is what a later \label writes down: the value \ref will answer with.
%
% SUBSTITUTION: LaTeX SNAPSHOTS the representation here -- \protected@edef
% freezes `2.1' into \@currentlabel, so a \label written after the next step of
% the same counter still records the number the step it belongs to produced.
% Nothing can be frozen here (see the four refusals above), so \@currentlabel is
% defined to READ the counter instead. A \label immediately after its
% \refstepcounter -- which is where LaTeX puts every one of them, inside the
% sectioning and caption commands -- records the same value either way. A \label
% separated from its own step by another step of the same counter does not.
\def\refstepcounter#1{\stepcounter{#1}\@refstepcounter@{#1}}
\def\@refstepcounter@#1{\def\@currentlabel{\csname the#1\endcsname}}
\def\@currentlabel{}
% latex.ltx: \def\value#1{\csname c@#1\endcsname}
\def\value#1{\expandafter\@valuectr\csname cr@#1\endcsname\@nil}
\def\@valuectr#1\@nil{\count#1}
% latex.ltx: the five representations, each applied to the counter's register.
\def\arabic#1{\expandafter\@arabicctr\csname cr@#1\endcsname\@nil}
\def\roman#1{\expandafter\@romanctr\csname cr@#1\endcsname\@nil}
\def\Roman#1{\expandafter\@Romanctr\csname cr@#1\endcsname\@nil}
\def\alph#1{\expandafter\@alphctr\csname cr@#1\endcsname\@nil}
\def\Alph#1{\expandafter\@Alphctr\csname cr@#1\endcsname\@nil}
\def\fnsymbol#1{\expandafter\@fnsymbolctr\csname cr@#1\endcsname\@nil}
\def\@arabicctr#1\@nil{\@arabic{\count#1}}
\def\@romanctr#1\@nil{\@roman{\count#1}}
\def\@Romanctr#1\@nil{\@Roman{\count#1}}
\def\@alphctr#1\@nil{\@alph{\count#1}}
\def\@Alphctr#1\@nil{\@Alph{\count#1}}
\def\@fnsymbolctr#1\@nil{\@fnsymbol{\count#1}}
% latex.ltx: \def\@arabic#1{\number #1}
\def\@arabic#1{\number #1}
% latex.ltx: \def\@alph#1{\ifcase#1\or a\or b\or ...\else\@ctrerr\fi}
\def\@alph#1{%
  \ifcase#1\or a\or b\or c\or d\or e\or f\or g\or h\or i\or j\or k\or l\or m\or
    n\or o\or p\or q\or r\or s\or t\or u\or v\or w\or x\or y\or z\else\@ctrerr\fi}
\def\@Alph#1{%
  \ifcase#1\or A\or B\or C\or D\or E\or F\or G\or H\or I\or J\or K\or L\or M\or
    N\or O\or P\or Q\or R\or S\or T\or U\or V\or W\or X\or Y\or Z\else\@ctrerr\fi}
% latex.ltx: \def\@fnsymbol#1{\ifcase#1\or *\or \dagger\or \ddagger\or ...}
% The symbols are the characters rather than the math atoms LaTeX sets them as,
% for the same reason \texttt yields text here.
\def\@fnsymbol#1{%
  \ifcase#1\or *\or \dagger\or \ddagger\or \S\or \P\or \|\or **\or
    \dagger\dagger\or \ddagger\ddagger\else\@ctrerr\fi}
\def\@ctrerr{??}
% SUBSTITUTION. latex.ltx writes \@roman with the primitive:
%
%     \def\@roman#1{\romannumeral#1}
%
% texrs has no \romannumeral. What is written out instead is the conversion
% tex.web S69 performs -- walk the value/letters table from the top, emit the
% letters and subtract while the remainder is at least the value -- with the
% loop UNROLLED to its bound. tex.web recurses; a recursive macro that takes
% arguments cannot be lowered here, because the lowerer inlines a macro's body
% and a body holding its own call is inlined forever. A letter repeats at most
% three times in a roman numeral, so three copies of each repeatable step IS the
% whole loop. The remainder is \count254 rather than a macro, for the reason the
% counters themselves are registers.
%
% The bound is 3999, which is where roman numerals stop having a spelling made
% of these letters. Above it \romannumeral goes on emitting m's; this says
% \@ctrerr, which is what LaTeX's own \@alph does when asked for a letter it has
% not got.
%
% WHERE IT DOES NOT WORK, and why it is the only representation with a place
% like this: the scratch register makes \@roman an ASSIGNMENT, and an assignment
% is not expandable, so anywhere a macro is expanded WITHOUT being executed the
% characters `\count 254=\count 100 \relax' come out instead of a numeral.
% \@arabic (`\number #1') and \@alph (`\ifcase#1...') have no such place: both
% are read by the message walker as a run-time read of the register.
% Measured: a \label inside a \part -- \thepart is \Roman{part} -- wrote
% `\newlabel{p:a}{{\count 254=\count 100 \relax }{0}}' into the .aux, which
% `src/latex/aux.rs' now drops rather than carry into the next run.
% The assignment-free form is a nest of \ifnum over absolute thresholds, and it
% is four thousand branches deep, because the remainder after each place is a
% subtraction and there is nowhere but a register to put one.
\def\@roman#1{\count254=#1 \relax\@rmn@guard\@rmn@list}
\def\@Roman#1{\count254=#1 \relax\@rmn@guard\@rmn@LIST}
\def\@rmn@guard{\ifnum\count254>3999 \@ctrerr\count254=0 \fi}
\def\@rmn@d#1#2{\ifnum\count254<#1 \else#2\advance\count254 by -#1 \fi}
\def\@rmn@list{%
  \@rmn@d{1000}{m}\@rmn@d{1000}{m}\@rmn@d{1000}{m}%
  \@rmn@d{900}{cm}\@rmn@d{500}{d}\@rmn@d{400}{cd}%
  \@rmn@d{100}{c}\@rmn@d{100}{c}\@rmn@d{100}{c}%
  \@rmn@d{90}{xc}\@rmn@d{50}{l}\@rmn@d{40}{xl}%
  \@rmn@d{10}{x}\@rmn@d{10}{x}\@rmn@d{10}{x}%
  \@rmn@d{9}{ix}\@rmn@d{5}{v}\@rmn@d{4}{iv}%
  \@rmn@d{1}{i}\@rmn@d{1}{i}\@rmn@d{1}{i}}
% The capital form. LaTeX reaches it through \@slowromancap over the lowercase
% letters; \uppercase is not a primitive here either, so the table is written
% again in capitals rather than a case change being faked.
\def\@rmn@D#1#2{\ifnum\count254<#1 \else#2\advance\count254 by -#1 \fi}
\def\@rmn@LIST{%
  \@rmn@D{1000}{M}\@rmn@D{1000}{M}\@rmn@D{1000}{M}%
  \@rmn@D{900}{CM}\@rmn@D{500}{D}\@rmn@D{400}{CD}%
  \@rmn@D{100}{C}\@rmn@D{100}{C}\@rmn@D{100}{C}%
  \@rmn@D{90}{XC}\@rmn@D{50}{L}\@rmn@D{40}{XL}%
  \@rmn@D{10}{X}\@rmn@D{10}{X}\@rmn@D{10}{X}%
  \@rmn@D{9}{IX}\@rmn@D{5}{V}\@rmn@D{4}{IV}%
  \@rmn@D{1}{I}\@rmn@D{1}{I}\@rmn@D{1}{I}}
%
% ---- the counters a class declares (article.cls / report.cls) -------------
% article.cls: \newcounter{part} \newcounter{section} \newcounter{subsection}[section]
% and so on; report.cls adds \newcounter{chapter} above them, and its
% \thesection is `\thechapter.\arabic{section}' where article's is
% `\arabic{section}'. report's is used, because the corpus sets books.
%
% Their registers, written out. A counter a DOCUMENT declares gets its register
% from `src/latex/counters.rs', which reads the source for \newcounter and
% \newtheorem and allocates from 120 upward; these are the ones every LaTeX
% document has whether it declares them or not, so they are fixed here and the
% Rust side starts past them.
%
% \@ctrreg{NAME}{NN} is the one place a register number is written: it makes
% both names that carry it -- \cr@NAME, the digits the READS expand through, and
% \ctr@NAME, the WRITE. `src/latex/counters.rs' emits the same two names for a
% counter the document declares, in the same shape.
\def\@ctrreg#1#2{\@namedef{cr@#1}{#2}\@namedef{ctr@#1}##1##2{##1{#2}{##2}}}
\@ctrreg{part}{100}
\@ctrreg{chapter}{101}
\@ctrreg{section}{102}
\@ctrreg{subsection}{103}
\@ctrreg{subsubsection}{104}
\@ctrreg{paragraph}{105}
\@ctrreg{subparagraph}{106}
\@ctrreg{figure}{107}
\@ctrreg{table}{108}
\@ctrreg{equation}{109}
\@ctrreg{footnote}{110}
\@ctrreg{mpfootnote}{111}
\@ctrreg{page}{112}
\@ctrreg{secnumdepth}{113}
\@ctrreg{tocdepth}{114}
\@ctrreg{enumi}{115}
\@ctrreg{enumii}{116}
\@ctrreg{enumiii}{117}
\@ctrreg{enumiv}{118}
\@ctrreg{Hfootnote}{119}
\newcounter{part}
\newcounter{chapter}
\newcounter{section}[chapter]
\newcounter{subsection}[section]
\newcounter{subsubsection}[subsection]
\newcounter{paragraph}[subsubsection]
\newcounter{subparagraph}[paragraph]
\newcounter{figure}[chapter]
\newcounter{table}[chapter]
\newcounter{equation}[chapter]
\newcounter{footnote}[chapter]
\newcounter{mpfootnote}
\newcounter{page}
\newcounter{secnumdepth}
\newcounter{tocdepth}
\newcounter{enumi}
\newcounter{enumii}
\newcounter{enumiii}
\newcounter{enumiv}
\newcounter{Hfootnote}
% The three counters LaTeX starts at something other than nought. WRITTEN, AND
% NOT IN EFFECT: this file is read through `Lowerer::preload', which keeps the
% definitions a preamble makes and DROPS the commands it emits, and a counter is
% a register, so an assignment made here is one of the dropped ones. Measured:
% `\documentclass{article}' then `\message{[\arabic{secnumdepth}]}' answers
% `[0]'. They are left here because they are what latex.ltx sets and because the
% moment a preamble's commands survive they will be right; what depends on them
% -- \pageref answering 0 rather than 1, and the secnumdepth test \@sect@run does
% not make -- says so where it depends on them.
\setcounter{page}{1}
\setcounter{secnumdepth}{2}
\setcounter{tocdepth}{2}
% report.cls's own \the<counter> definitions.
\def\thepart{\Roman{part}}
\def\thechapter{\arabic{chapter}}
\def\thesection{\thechapter.\arabic{section}}
\def\thesubsection{\thesection.\arabic{subsection}}
\def\thesubsubsection{\thesubsection.\arabic{subsubsection}}
\def\theparagraph{\thesubsubsection.\arabic{paragraph}}
\def\thesubparagraph{\theparagraph.\arabic{subparagraph}}
\def\thefigure{\thechapter.\arabic{figure}}
\def\thetable{\thechapter.\arabic{table}}
\def\theequation{\thechapter.\arabic{equation}}
\def\thefootnote{\arabic{footnote}}
\def\thepage{\arabic{page}}
\def\theenumi{\arabic{enumi}}
\def\theenumii{\alph{enumii}}
\def\theenumiii{\roman{enumiii}}
\def\theenumiv{\Alph{enumiv}}
%
% ---- cross references: ltxref.dtx ----------------------------------------
% latex.ltx:
%   \def\newlabel{\@newl@bel r}
%   \def\@newl@bel#1#2#3{...\global\@namedef{#1@#2}{#3}}
% An .aux file is a list of these. texrs reads the .aux in Rust
% (`src/latex/aux.rs') and hands the entries to a run as \newlabel calls in the
% preamble, which is the same round trip LaTeX makes through \@input.
\def\newlabel#1#2{\@namedef{r@#1}{#2}}
\def\bibcite#1#2{\@namedef{b@#1}{#2}}
% latex.ltx:
%   \def\ref#1{\expandafter\@setref\csname r@#1\endcsname\@firstoftwo{#1}}
%   \def\pageref#1{\expandafter\@setref\csname r@#1\endcsname\@secondoftwo{#1}}
%   \def\@setref#1#2#3{\ifx#1\relax ...??... \else\expandafter#2#1\null\fi}
%
% SUBSTITUTION in \@setref only. `\ifx#1\relax' is how LaTeX asks whether the
% label was ever written down, and in texrs that test answers FALSE for an
% undefined control sequence -- `\expandafter\ifx\csname r@x\endcsname\relax'
% takes the \else arm for a name nothing has defined, and \ifcsname is refused
% outright with `! Unsupported conditional \ifcsname.' So the question is asked
% where it CAN be answered: `src/latex/aux.rs' knows every label the .aux
% defines and every label the document references, and defines \r@X as
% `{??}{??}' for the ones referenced and not defined. \@setref then never meets
% an undefined name, and `??' -- LaTeX's own answer for an unresolved reference
% -- comes out of the same place it comes out of under LaTeX.
% \ref, \pageref and \label are NOT defined here: the lowerer claims them by
% name (lower_cross_ref) and marks the place, and the pass that lays out the
% pages answers them -- in ONE run, where this two-pass form answers `??' until
% a second. \@setref stays for a class that builds its own.


\def\@setref#1#2#3{\expandafter#2#1}
% amsmath: \eqref{...} is \ref in parentheses, through \tagform@.
\def\eqref#1{\textup{\tagform@{\ref{#1}}}}
\def\tagform@#1{(#1)}
\def\textup#1{#1}
% latex.ltx:
%   \def\label#1{...\protected@write\@auxout{}%
%      {\string\newlabel{#1}{{\@currentlabel}{\thepage}}}...}
% There is no \write here -- nor \openin, nor \read -- so the writing side is
% Rust's: `src/latex/aux.rs' runs the document a second time with \label \let to
% \@auxlabel, which puts the entry in the \message stream, and writes the .aux
% from what came back. In an ordinary run \label produces nothing, exactly as it
% does under LaTeX: a label is a target, not text.

\def\@auxlabel#1{\message{\string\newlabel{#1}{{\@currentlabel}{\thepage}}}}
%
% ---- sectioning: ltsect.dtx ----------------------------------------------
% latex.ltx:
%   \def\@sect#1#2#3#4#5#6[#7]#8{%
%     \ifnum #2>\c@secnumdepth \let\@svsec\@empty
%     \else \refstepcounter{#1}%
%           \protected@edef\@svsec{\@seccntformat{#1}\relax}\fi ...}
%   \def\@startsection#1#2#3#4#5#6{... \@ifstar{\@ssect...}{\@dblarg{\@sect...}}}
%
% What is carried over is the two things a sectioning command does that a
% mouth can do: it STEPS its counter, and it yields its title. What is dropped
% is everything \@startsection's six arguments describe -- the space above and
% below, the indent, the font -- because that is the stomach's, and in a run
% that is setting text `src/lower.rs' claims \section, \subsection,
% \subsubsection and \chapter before this definition is reached and gives them
% the breaks and the heading marks a page needs.
%
% Stepping the counter is what makes a \label inside a section record the
% SECTION's number: \refstepcounter is where \@currentlabel is set, and without
% it every entry written into the .aux carried an empty value.
%
% SUBSTITUTION: latex.ltx's `\ifnum #2>\c@secnumdepth' is dropped and the
% counter is always stepped. \c@secnumdepth would have to be right to test it,
% and it is not: the preamble's `\setcounter{secnumdepth}{2}' is read by
% `Lowerer::preload', which keeps the DEFINITIONS a preamble makes and drops
% the commands, so every counter starts a run at zero whatever the preamble set
% it to. Measured: `\documentclass{article}' then
% `\message{[\arabic{secnumdepth}]}' answers `[0]'. Stepping regardless numbers
% a section the document asked not to number; testing against nought would
% number nothing at all.
%
% The optional argument is the entry a table of contents would have carried,
% and there is no contents in the mouth, so it is read and dropped.
\def\@sect@run#1#2{\refstepcounter{#1}#2}
\def\@sect@opt#1[#2]#3{\@sect@run{#1}{#3}}
\def\@sect@plain#1{\@ifnextchar[{\@sect@opt{#1}}{\@sect@run{#1}}}
% The starred form takes no number, so no counter is stepped: latex.ltx sends it
% to \@ssect, which is \@sect without the \refstepcounter.
\def\@sect@#1{\@ifstar\@firstofone{\@sect@plain{#1}}}
%
% ---- the bibliography: ltbibl.dtx -----------------------------------------
% latex.ltx:
%   \def\@biblabel#1{[#1]}
%   \def\bibitem{\@ifnextchar[\@lbibitem\@bibitem}
%   \def\@lbibitem[#1]#2{\item[\@biblabel{#1}\hfill]...\bibcite{#2}{#1}...}
%   \def\@bibitem#1{\item ...\bibcite{#1}{\the\value{\@listctr}}...\ignorespaces}
%   \def\@cite#1#2{[{#1\if@tempswa , #2\fi}]}
%
% \bibitem was UNDEFINED here, and an undefined control sequence stops a
% document: `\begin{thebibliography}{9}\bibitem{a} An entry.' ended the run at
% `! Undefined control sequence \bibitem.' and lost everything after it.
%
% The number a \bibitem gets is its position in the list, which latex.ltx reads
% off the enclosing list's counter; enumiv is the counter a \bibitem list uses
% and it is stepped here directly, because there is no list to ask. The label is
% \@biblabel's, so an entry reads `[1] An entry.' as it does under LaTeX --
% minus the hanging indent, which is the stomach's.
%
% \cite answers out of \b@KEY, which is what \bibcite defines, exactly as \ref
% answers out of \r@KEY: the .aux carries both and `src/latex/aux.rs' seeds a
% key nothing defined -- with `?', which is LaTeX's own answer for a citation
% with no entry.
\def\@biblabel#1{[#1]}
% latex.ltx's list opens `\usecounter{enumiv}' with `\def\@listctr{enumiv}' and
% `\let\p@enumiv\@empty', and it is where \theenumiv stops being the enumerate
% form (\Alph) and becomes the entry number.
\def\thebibliography#1{\setcounter{enumiv}{0}\def\theenumiv{\arabic{enumiv}}}
\def\endthebibliography{}
\def\bibitem{\@ifnextchar[\@lbibitem\@bibitem}
\def\@lbibitem[#1]#2{\@biblabel{#1} }
% The label printed is \b@KEY -- the number the .aux carries -- and NOT
% \@ctrtotext{enumiv}, though the counter is stepped for the pass that writes
% that file. The reason is the one \@ctrtotext exists for: a counter cannot be
% read into the document's text, so an entry labelled from the counter came out
% as `[]'. \b@KEY is characters, seeded by `src/latex/aux.rs', and characters
% reach the text. It is also LaTeX's own data flow -- the number a \cite prints
% comes from the .aux either way -- so a first run says `[?]' and the run after
% it says `[1]', which is what latex does.
\def\@bibitem#1{\stepcounter{enumiv}\@biblabel{\@nameuse{b@#1}} }
% The diversion `src/latex/aux.rs' runs the document under to collect the
% entries, beside \@auxlabel and for the same reason.
\def\@auxbibitem{\@ifnextchar[\@auxlbibitem\@auxbibitem@}
\def\@auxlbibitem[#1]#2{\message{\string\bibcite{#2}{#1}}}
\def\@auxbibitem@#1{\stepcounter{enumiv}%
  \message{\string\bibcite{#1}{\@ctrtotext{enumiv}}}}
\def\cite{\@ifnextchar[\@citeopt\@cite@}
\def\@citeopt[#1]#2{\@cite@{#2}}
\def\@cite@#1{\@biblabel{\@nameuse{b@#1}}}
\def\nocite#1{}
\def\part{\@sect@{part}}
\def\chapter{\@sect@{chapter}}
\def\section{\@sect@{section}}
\def\subsection{\@sect@{subsection}}
\def\subsubsection{\@sect@{subsubsection}}
\def\paragraph{\@sect@{paragraph}}
\def\subparagraph{\@sect@{subparagraph}}
%
% ---- \newenvironment: ltdefns.dtx ----------------------------------------
% latex.ltx:
%   \def\newenvironment{\@star@or@long\new@environment}
%   \def\new@environment#1{\@ifundefined{#1}{...}...\@newenv{#1}}
%   \def\@newenv#1#2#3#4{\@namedef{#1}...\@namedef{end#1}{#4}}
% \begin{x} runs \x and \end{x} runs \endx (latex.ltx, ltmiscen), which the
% prelude already does. What was missing is that \newenvironment threw both
% bodies away, so a document that defined its own environment got neither of
% them. The argument count and the optional default are handed to \newcommand,
% which is what latex.ltx does too.
%
% The \@ifundefined guard latex.ltx opens with is dropped rather than faked: it
% exists to REFUSE a redefinition, and refusing needs the test that does not
% work here. A document that defines the same environment twice gets the second
% definition instead of an error.
\def\newenvironment{\@ifstar\@newenv@\@newenv@}
\def\renewenvironment{\@ifstar\@newenv@\@newenv@}
\def\@newenv@#1{\@ifnextchar[{\@newenv@opt{#1}}{\@newenv@plain{#1}}}
\long\def\@newenv@plain#1#2#3{%
  \expandafter\def\csname #1\endcsname{#2}%
  \expandafter\def\csname end#1\endcsname{#3}}
\def\@newenv@opt#1[#2]{\@ifnextchar[{\@newenv@optdef{#1}{#2}}{\@newenv@args{#1}{#2}}}
\long\def\@newenv@args#1#2#3#4{%
  \expandafter\newcommand\csname #1\endcsname[#2]{#3}%
  \expandafter\def\csname end#1\endcsname{#4}}
\long\def\@newenv@optdef#1#2[#3]#4#5{%
  \expandafter\newcommand\csname #1\endcsname[#2][#3]{#4}%
  \expandafter\def\csname end#1\endcsname{#5}}
% latex.ltx: \def\@ifdefinable#1#2{...}, which \newcommand uses to refuse a name
% that is already taken. The refusal needs \ifx against \relax; the
% definable-name half is what callers here want, so that is what is kept.
\def\@ifdefinable#1#2{#2}
%
% ---- a counter on its way into the DOCUMENT'S TEXT -----------------------
% Every read of a counter that ends up as WORDS goes through \@ctrtotext, and
% there is one reason for it: \number and \the do not reach the text stream.
% `\count110=7 A\number\count110 B' under --text answers
% `! Undefined control sequence \number.', and \the answers the same; both work
% inside \message, where the lowerer keeps them as a run-time read, and while
% \edef scans an operand. So \arabic -- which latex.ltx defines as `\number #1'
% and which this file ports unchanged -- is usable everywhere a counter is
% TESTED or MESSAGED and nowhere a counter is PRINTED.
%
% `src/latex.rs' redefines this to produce nothing when the run is setting text,
% which is what a footnote mark and a caption number produced before this file
% existed, so nothing regresses. The moment the lowerer can splice a register
% into text the way it already splices one into a \message, the redefinition
% goes and the numbers appear.
\def\@ctrtotext#1{\csname the#1\endcsname}
%
% ---- footnotes: ltfloat.dtx ----------------------------------------------
% latex.ltx:
%   \def\footnote{\@ifnextchar[\@xfootnote\@yfootnote}
%   \def\@yfootnote{\stepcounter\@mpfn
%      \protected@xdef\@thefnmark{\thempfn}\@footnotemark\@footnotetext}
%   \def\@xfootnote[#1]{\begingroup\csname c@\@mpfn\endcsname #1\relax
%      \unrestored@protected@xdef\@thefnmark{\thempfn}\endgroup
%      \@footnotemark\@footnotetext}
%
% The numbering is LaTeX's: the counter steps, \@thefnmark is its
% representation, and that mark is what \footnotemark alone also produces. Where
% the note GOES is the difference -- LaTeX puts it in an \insert and the output
% routine sets it at the foot of the page, and there is no \insert here, so the
% note follows its mark in the text. That is the reading the rest of this layer
% takes: the words survive, the page does not.
%
% \@thefnmark is \def rather than \protected@xdef for the reason \@currentlabel
% is: nothing freezes here.
\def\@mpfn{footnote}
\def\thempfn{\@ctrtotext{footnote}}
\def\footnote{\@ifnextchar[\@xfootnote\@yfootnote}
\long\def\@yfootnote#1{%
  \stepcounter{footnote}\def\@thefnmark{\thempfn}%
  \@footnotemark\@footnotetext{#1}}
\long\def\@xfootnote[#1]#2{%
  \setcounter{footnote}{#1}\def\@thefnmark{\thempfn}%
  \@footnotemark\@footnotetext{#2}}
\def\@footnotemark{\textsuperscript{\@thefnmark}}
\long\def\@footnotetext#1{ [\@thefnmark] #1 }
\def\footnotemark{\@ifnextchar[\@xfootnotemark\@yfootnotemark}
\def\@yfootnotemark{\stepcounter{footnote}\def\@thefnmark{\thempfn}\@footnotemark}
\def\@xfootnotemark[#1]{\setcounter{footnote}{#1}\def\@thefnmark{\thempfn}\@footnotemark}
\def\footnotetext{\@ifnextchar[\@xfootnotetext\@yfootnotetext}
\long\def\@yfootnotetext#1{\@footnotetext{#1}}
\long\def\@xfootnotetext[#1]#2{%
  \setcounter{footnote}{#1}\def\@thefnmark{\thempfn}\@footnotetext{#2}}
%
% ---- captions: ltfloat.dtx -----------------------------------------------
% latex.ltx:
%   \def\caption{\refstepcounter\@captype \@dblarg{\@caption\@captype}}
%   \long\def\@caption#1[#2]#3{...\@makecaption{\csname fnum@#1\endcsname}%
%      {\ignorespaces #3}...}
%   \def\fnum@figure{\figurename\nobreakspace\thefigure}
%
% \@captype is set by the float environment -- `figure' inside a figure,
% `table' inside a table. The number is the counter's, stepped by
% \refstepcounter, which is also what makes a \label inside a caption record the
% FIGURE's number rather than the section's. What LaTeX does with the text --
% set it in a box of the float's width, centred if it fits on one line -- is the
% stomach's, so the caption reads as its own run of text here.
\def\@captype{figure}
\def\figurename{Figure}
\def\tablename{Table}
\def\fnum@figure{\figurename\ \@ctrtotext{figure}}
\def\fnum@table{\tablename\ \@ctrtotext{table}}
\def\caption{\refstepcounter\@captype\@ifnextchar[{\@caption\@captype}{\@caption\@captype[]}}
\long\def\@caption#1[#2]#3{\csname fnum@#1\endcsname: #3}
%
% ---- the option and file machinery: ltclass.dtx ---------------------------
% What `\usepackage{x}' does when the package really loads. The loading itself
% is Rust's -- `src/latex/load.rs' finds x.sty with kpsewhich and puts it in the
% preamble -- and these are the declarations the .sty file's own text needs in
% order to be read at all.
%
% latex.ltx:18801-18821, reduced to the classic form. The 2025 kernel routes
% \@pass@ptions through the expl3 file hooks
% (\@expl@@@filehook@set@curr@file@@nNN, \tl_trim_spaces:n); no expl3 runs here,
% so what is carried is the option-list append those hooks are wrapped around.
\def\@ptionlist#1{\@nameuse{opt@#1}}
\def\@pass@ptions#1#2#3{\@namedef{opt@#3.#1}{#2}}
\def\@pkgextension{sty}
\def\@clsextension{cls}
% \PassOptionsToPackage and \PassOptionsToClass are NOT defined here, and the
% omission is deliberate. The lowerer reads both by name -- `src/lower.rs',
% "PassOptionsToPackage" | "PassOptionsToClass" -- because they are where a
% pandoc preamble states the page it is set on:
%
%     \PassOptionsToPackage{margin=0.5in}{geometry}
%
% is the only place that half-inch is written, and `Layout::absorb_geometry_
% options' takes it from there. That dispatch is by name and only for a name
% nothing has defined as a macro, so defining either of them here TOOK the
% options away from the page: `tests/typeset.rs',
% options_passed_to_the_class_and_to_geometry_reach_the_page_too, set its text
% at 72pt rather than at the 36pt it asked for.
%
% What is lost is the recording -- \@ptionlist{x.sty} answers nothing for a
% package the preamble passed options to -- and the lowerer consumes both
% arguments, so a .sty that calls either one still reads to the end.
% latex.ltx:18824-18831.
\def\DeclareOption{\@ifstar\@defdefault@ds\@declareoption}
\long\def\@declareoption#1#2{%
  \xdef\@declaredoptions{\@declaredoptions,#1}%
  \@namedef{ds@#1}{#2}}
\long\def\@defdefault@ds#1{\def\default@ds{#1}}
\def\@declaredoptions{}
\def\default@ds{}
\def\CurrentOption{}
% SUBSTITUTION: the comma list is walked UNROLLED, not with \@for.
%
% latex.ltx spells both \ProcessOptions and \ExecuteOptions as \@for over the
% option list, and \@for is carried above verbatim -- and does not run here.
% Measured: `\@for\x:=a,b\do{\message{<\x>}}' answers `! Extra \else.', because
% \@iforloop calls ITSELF and the lowerer inlines a macro's body, so a body
% holding its own call cannot be lowered (`src/latex/kernel.tex' says the same
% thing about tex.web's \romannumeral). Every class stopped on the line that
% used it: article.cls:111 is \ExecuteOptions{letterpaper,10pt,oneside,...}, and
% `texrs: package article is not loadable: Extra \fi' was that.
%
% What replaces it is the same unrolling \@roman uses: a delimited parameter
% text takes eight items at a time, three of them in a row, and the list is
% padded with \@nil so a shorter list still matches. Twenty-four options is the
% bound -- article.cls's own \ExecuteOptions has five, and a \documentclass line
% rarely has three -- and an option past it is dropped rather than run.
%
% \@opt@one is the loop body: latex.ltx's `\ifx\CurrentOption\@empty\else...',
% with the \@nil guard the padding needs beside it.
\def\@opt@one#1{\def\CurrentOption{#1}%
  \ifx\CurrentOption\@nnil\else\ifx\CurrentOption\@empty\else
    \@nameuse{ds@#1}\fi\fi}
\def\@opt@eight#1#2#3#4#5#6#7#8{%
  \@opt@one{#1}\@opt@one{#2}\@opt@one{#3}\@opt@one{#4}%
  \@opt@one{#5}\@opt@one{#6}\@opt@one{#7}\@opt@one{#8}}
\def\@opt@A#1,#2,#3,#4,#5,#6,#7,#8,#9\@@{%
  \@opt@eight{#1}{#2}{#3}{#4}{#5}{#6}{#7}{#8}\@opt@B#9\@@}
\def\@opt@B#1,#2,#3,#4,#5,#6,#7,#8,#9\@@{%
  \@opt@eight{#1}{#2}{#3}{#4}{#5}{#6}{#7}{#8}\@opt@C#9\@@}
% The last eight, and #9 -- everything past twenty-four -- dropped.
\def\@opt@C#1,#2,#3,#4,#5,#6,#7,#8,#9\@@{%
  \@opt@eight{#1}{#2}{#3}{#4}{#5}{#6}{#7}{#8}}
% Twenty-five \@nil, so a list of k options leaves the last level a ninth
% parameter to swallow for every k from 0 to 24.
\def\@do@ptions#1{\@opt@A#1,\@nil,\@nil,\@nil,\@nil,\@nil,\@nil,\@nil,\@nil,%
  \@nil,\@nil,\@nil,\@nil,\@nil,\@nil,\@nil,\@nil,\@nil,\@nil,\@nil,\@nil,%
  \@nil,\@nil,\@nil,\@nil,\@nil\@@\let\CurrentOption\@empty}
% latex.ltx:18847-18899, reduced: \ProcessOptions runs the declared option code
% for every option the document asked for. The unused-option bookkeeping and the
% \detokenize round trip are dropped -- \detokenize is not a primitive here.
% \@curroptions is set by `src/latex/load.rs' from the [...] the document wrote,
% which is where \usepackage would have taken it from.
\def\ProcessOptions{\@ifstar\@process@ptions\@process@ptions}
\def\@process@ptions{\expandafter\@do@ptions\expandafter{\@curroptions}}
\def\@curroptions{}
\def\@options{\ProcessOptions*}
% latex.ltx:18901-18907.
\def\ExecuteOptions#1{\@do@ptions{#1}}
% latex.ltx:18955-18974. texrs is not the LaTeX2e format, and saying so would
% stop every .sty file on its first line; what the test guards is a version
% warning, so it is answered the way a matching format answers it.
\def\fmtname{LaTeX2e}
\def\fmtversion{2025/11/01}
\def\NeedsTeXFormat#1{\@ifnextchar[\@needsf@rmat@gobble\@empty}
\def\@needsf@rmat@gobble[#1]{}
% latex.ltx:18763-18789, reduced to what it records: the version string, under
% the name every \@ifpackageloaded test reads.
\def\ProvidesPackage#1{\@ifnextchar[{\@pr@videpackage{#1}}{\@pr@videpackage{#1}[]}}
\def\@pr@videpackage#1[#2]{\@namedef{ver@#1.\@currext}{#2}}
\def\ProvidesClass#1{\@ifnextchar[{\@pr@videclass{#1}}{\@pr@videclass{#1}[]}}
\def\@pr@videclass#1[#2]{\@namedef{ver@#1.cls}{#2}}
\def\ProvidesFile#1{\@ifnextchar[\@providesfile@gobble\@empty}
\def\@providesfile@gobble[#1]{}
\def\@currext{sty}
\def\@currname{}
% latex.ltx:18697-18716. The date comparison every package opens with, and the
% single most-reported thing a package wanted: `texrs: package xcolor needs
% \@ifl@t@r' was the first line of the run for 54 of the 274 corpus documents.
%
%   \def\@ifl@t@r#1#2{\ifnum\expandafter\@parse@version@#1//00\@nil<%
%                           \expandafter\@parse@version@#2//00\@nil
%     \expandafter\@secondoftwo \else \expandafter\@firstoftwo \fi}
%
% Carried whole, delimiters and all: `2025/11/01' becomes the number 20251101
% through \@parse@version's parameter text, and the version is compared as an
% integer. \fmtversion is this file's own, above.
% The date PARSE is carried whole -- `2025/11/01' becomes 020251101 through
% these three parameter texts, and `\count200=\expandafter\@parse@version@
% \fmtversion//00\@nil' really holds 20251101.
%
% One \if of latex.ltx's is dropped: its \@parse@version@dash is
%
%     \def\@parse@version@dash#1-#2-#3#4#5\@nil{\if\relax#2\relax\else#1\fi#2#3#4 }
%
% and that \if sits inside the number \ifnum is scanning, which this engine's
% number scanner does not expand -- `! Missing = inserted for \ifnum.' It drops
% the year when the MONTH is empty, and the only version with an empty month is
% the empty one, where #1 is the bare `0' \@parse@version@ prepended: dropping
% it and keeping it both give nought.
\def\@parse@version@#1{\@parse@version0#1}
\def\@parse@version#1/#2/#3#4#5\@nil{\@parse@version@dash#1-#2-#3#4\@nil}
\def\@parse@version@dash#1-#2-#3#4#5\@nil{#1#2#3#4 }
% SUBSTITUTION, and the COMPARISON is what is substituted for. latex.ltx picks
% between \@firstoftwo and \@secondoftwo inside an \ifnum, and neither shape of
% that works here:
%
%   \expandafter\@secondoftwo\else\expandafter\@firstoftwo\fi  answers
%       `! Extra \else.' -- an arm expanded while the conditional is open leaves
%       the \fi for whatever the arm reads next, which `prelude.tex' explains at
%       \@ifnch;
%   \let\@arm\@secondoftwo\else\let\@arm\@firstoftwo\fi\@arm   assigns TWICE.
%       An \ifnum at the top level of a file is lowered as a run-time branch and
%       BOTH arms are lowered, so both \lets run and the second wins. Measured:
%       `\ifnum 1<2 \let\a\@secondoftwo\else\let\a\@firstoftwo\fi \a{X}{Y}' says
%       X, and so does the same line with `5<2'. An \ifnum is never decided
%       while lowering, whatever its operands are.
%
% So the answer given is the one this format's own version makes true.
% \fmtversion is 2025/11/01 and every use of this that reaches a package's first
% page is `\IfFormatAtLeastTF{2022/06/01}' or older -- xcolor.sty:48-51,
% ifthen.sty:75 -- so `later' is the right arm and the arm that goes on reading
% the file: xcolor's else arm is \endinput. \@ifpackagelater, which shares this,
% then says a loaded package is new enough, which is likewise the arm that keeps
% the document running.
\def\@ifl@t@r#1#2{\@firstoftwo}
\def\IfFormatAtLeastTF{\@ifl@t@r\fmtversion}
\def\@ifl@ter#1#2{\expandafter\@ifl@t@r\csname ver@#2.#1\endcsname}
\def\@ifpackagelater{\@ifl@ter\@pkgextension}
\def\@ifclasslater{\@ifl@ter\@clsextension}
\let\IfPackageAtLeastTF\@ifpackagelater
\let\IfClassAtLeastTF\@ifclasslater
% latex.ltx:19464 and its neighbours. The whole body of each is inside
% `\ifnum\pkgcls@targetdate>\z@' -- "some sort of rollback request" -- and
% \pkgcls@targetdate is zero unless the document asked for one with
% `\usepackage[=2020-02-02]{amsmath}'. Nothing here asks, so what these do IS
% nothing, and that is what is written rather than the rollback machinery they
% guard. They were the second most-reported name: amsmath, textcomp, parskip,
% longtable and array all open with one.
\def\DeclareRelease#1#2#3{}
\def\DeclareCurrentRelease#1#2{}
% latex.ltx:18687-18693. The test is `\ifx\csname ver@x.sty\endcsname\relax',
% which is the one texrs cannot answer, so it goes through \@ifundefined -- see
% the note on it directly below.
\def\@ifpackageloaded#1{\@ifundefined{ver@#1.sty}}
\def\@ifclassloaded#1{\@ifundefined{ver@#1.cls}}
% latex.ltx:1737. \@ifundefined#1{true}{false} is true when NOTHING is defined
% under that name:
%
%     \def\@ifundefined#1{\expandafter\ifx\csname #1\endcsname\relax
%       \expandafter\@firstoftwo\else\expandafter\@secondoftwo\fi}
%
% PORTED now, and this was the substitution that cost the most. The reading it
% replaces treated every name as DEFINED, because texrs answers that \ifx FALSE
% for a name nothing has defined -- an undefined control sequence is not
% \ifx-equal to \relax here -- so a .sty that branched on \@ifundefined to
% SUPPLY a missing macro did not supply it.
%
% \ifcsname is what asks the question, and the engine has it now
% (`src/expand.rs'). One shape change from latex.ltx's: the arms are chosen with
% \let inside the conditional and called after the \fi rather than reached
% through \expandafter, for the reason `prelude.tex' gives at \@ifnch -- an arm
% expanded while the conditional is still open leaves the \fi to whatever the
% arm reads next, and answers `! Extra \else.'
%
% The senses are latex.ltx's: the FIRST argument is what runs when nothing is
% defined under the name, so \ifcsname being true takes the second.
\def\@ifundefined#1{\ifcsname #1\endcsname
  \let\@ifu@arm\@secondoftwo\else\let\@ifu@arm\@firstoftwo\fi\@ifu@arm}
% latex.ltx (ltfiles). \IfFileExists asks the file system, which is Rust's job.
%
%     \DeclareRobustCommand\InputIfFileExists[2]{\IfFileExists{#1}{...\@@input}}
%
% WIRED, to `src/latex/load.rs': it writes one \let per file that really
% exists, out of the names this document and its packages name literally, and
% the sentinel below is what those \lets point at. The dispatch has to be a
% comparison rather than \@ifundefined because nothing here can ask whether a
% built name is defined -- see the note on \@ifundefined -- and measured, an
% \ifx against a sentinel IS decided while lowering and is false for a name
% nothing defined. So an unknown name, including one built by expansion, takes
% the not-found arm, which is the arm this macro took for every name before it
% was wired.
%
% The arms are chosen with \let inside the conditional and called after the
% \fi, for the reason `prelude.tex' gives at \@ifnch: an arm expanded while the
% conditional is still open leaves the \fi to whatever the arm reads next.
\def\@fe@yes{}
\def\IfFileExists#1#2#3{%
  \expandafter\ifx\csname @fe@#1\endcsname\@fe@yes
    \let\@fe@arm\@firstoftwo\else\let\@fe@arm\@secondoftwo\fi
  \@fe@arm{#2}{#3}}
% latex.ltx:9870. The found arm runs #2 and then reads the file, which is what
% \@@input does there; here it is the engine's own \input, whose braced form
% and TeX-tree search are what make a name like `textcomp.cfg' reachable at all.
\def\InputIfFileExists#1#2#3{\IfFileExists{#1}{#2\input{#1}}{#3}}
% latex.ltx (ltclass): the hooks.
%
% NOT PORTED, and this is the one place in this file where a stand-in is kept on
% purpose rather than for want of a primitive. \AtBeginDocument DEFERS its
% argument to \begin{document}, by appending it to \@begindocumenthook, and
% appending to a macro means reading the macro's body and writing it back --
% \g@addto@macro's `\xdef#1{\the\toks@}'. Nothing freezes in this engine, so
% the append cannot be written: every shape of it either loses the earlier
% entries or builds a macro that reads itself.
%
% PORTED now, and both halves of what blocked it have arrived. The append is
% \g@addto@macro, which this file writes with \expandafter rather than with
% \xdef (see its own note), and the thing the hooks in the corpus actually DO is
% ask \ifcsname -- scifi2/docs/book.tex registers
%
%     \AtBeginDocument{\ifcsname Shaded\endcsname\else...\fi ...}
%
% which refused outright until `src/expand.rs' had \ifcsname. Running the hook
% is what LaTeX does and is now what happens: \AtBeginDocument appends to
% \@begindocumenthook and \document, which is what \begin{document} expands,
% runs it.
%
% \AtEndDocument's hook is collected and NOT run, because nothing here reaches
% the end of the document as an event: \end{document} is the last thing read and
% there is no output routine after it to fire.
\def\@begindocumenthook{}
\def\@enddocumenthook{}
\def\AtBeginDocument#1{\g@addto@macro\@begindocumenthook{#1}}
\def\AtEndDocument#1{\g@addto@macro\@enddocumenthook{#1}}
\def\AtEndOfPackage#1{#1}
\def\AtEndOfClass#1{#1}
\def\@onlypreamble#1{}
% latex.ltx:1340. \newif, which nearly every package opens with:
%
%     \def\newif#1{\count@\escapechar \escapechar\m@ne
%       \let#1\iffalse \@if#1\@empty\iftrue \@if#1\@empty\iffalse
%       \escapechar\count@}
%     \def\@if#1#2{\expandafter\def\csname\expandafter\@gobbletwo\string#1%
%                    \expandafter\@gobbletwo\string#2\endcsname}
%
% HALF PORTED, and the missing half is said rather than faked. `\let#1\iffalse'
% is here and works: \newif\ifdraft makes \ifdraft a false conditional. What is
% NOT here is \ifdraft's pair of switches, \drafttrue and \draftfalse, because
% building their names means taking `ifdraft' apart and dropping the `if', and
% both ways TeX spells a control sequence as characters come back from this
% engine as ONE opaque token rather than as the characters:
%
%   \def\g if#1\@nil{...}\expandafter\g\csstring\ifdraft\@nil
%       answers `! Use of macro doesn't match its definition.', while the same
%       macro applied to the characters -- \g ifdraft\@nil -- matches and gives
%       `draft'. \string behaves the same way, and \escapechar, which is how
%       latex.ltx spells a name without its backslash, is not a parameter this
%       engine has at all.
%
% So a document that says \drafttrue stops on an undefined control sequence,
% loudly and on the line that wrote it, rather than silently taking the false
% branch of a switch it thought it had turned on.
\def\newif#1{\let#1\iffalse}
%
% ---- the document environment --------------------------------------------
% latex.ltx (ltmiscen): \begin{document} ends the preamble. There is no output
% routine to start here, so what it does is nothing -- but it has to BE defined,
% because \begin{document} expands \document and an undefined one is where a
% document with no \documentclass line would stop.
\def\document{\@begindocumenthook}
\def\enddocument{}
%
% ---- the format's own switches -------------------------------------------
% latex.ltx declares these with \newif, and a CLASS reads them without
% declaring them: article.cls opens with `\if@compatibility' and stops on the
% first line if the format has not made it. They are the format's, so they are
% ported here, and `src/latex/switches.rs` -- which is handed this file too --
% writes each one's \NAMEtrue and \NAMEfalse.
%
% The list is every \newif in
% /usr/local/texlive/2026/texmf-dist/tex/latex/base/latex.ltx, in its order,
% minus the four that are internal to code this file does not carry
% (\ifin@ is defined above, \if@includeinrelease guards a release wrapper the
% port has no use for, and \ifh@ and \ifv@ belong to the picture environment).
\newif\if@afterindent
\newif\if@compatibility
\newif\if@endpe
\newif\if@eqnsw
\newif\if@filesw
\newif\if@firstamp
\newif\if@font@series@context
\newif\if@forced@series
\newif\if@in@minipage@env
\newif\if@inlabel
\newif\if@negarg
\newif\if@newlist
\newif\if@nmbrlist
\newif\if@noitemarg
\newif\if@noparitem
\newif\if@noparlist
\newif\if@noskipsec
\newif\if@ovb
\newif\if@ovhline
\newif\if@ovl
\newif\if@ovr
\newif\if@ovt
\newif\if@ovvline
\newif\if@partsw
\newif\if@pboxsw
\newif\if@rjfield
\newif\if@tempswa
\newif\ifdt@p
\newif\ifmath@fonts
\newif\ifmaybe@ic
% latex.ltx:20674-20681, the output routine's own, written `\newif \if@insert'
% with a space. They are the ones a CLASS turns on and off without declaring:
% article.cls's first page of code says \@mparswitchfalse, and every document
% that loaded article stopped there --
% `texrs: package article needs \@mparswitchfalse' -- because the switch the
% format was supposed to have made was not there for switches.rs to pair.
\newif\if@insert
\newif\if@fcolmade
\newif\if@firstcolumn
\newif\if@reversemargin
\newif\if@mparswitch
% latex.ltx turns three of them on or off on the line that declares them, and
% \if@firstcolumn is the one whose initial value is not \newif's own false. It
% is set with \let rather than with \@firstcolumntrue because THIS file is what
% switches.rs reads to write \@firstcolumntrue: the pair does not exist yet on
% the line below the \newif that asks for it.
\let\if@firstcolumn\iftrue
% classes.dtx declares these; every standard class reads all of them, and a
% class that declares one itself simply gets it twice, which \newif allows.
\newif\if@titlepage
\newif\if@twoside
\newif\if@twocolumn
\newif\if@openright
\newif\if@mainmatter
\newif\if@restonecol
\newif\if@specialpage
\newif\if@dblfloat
\newif\if@nobreak
\newif\if@insertfootnotes
%
% ---- the dimensions a counter can be set from ----------------------------
% plain.tex:  \newdimen\maxdimen  \maxdimen=16383.99998pt
%
% The prelude answers \maxdimen with nothing, which is right for a run that
% cannot measure anything -- and wrong the moment a counter is really set,
% because pandoc's preamble opens with
%
%     \setcounter{secnumdepth}{-\maxdimen}   % remove section numbering
%
% and `\count113=-' is `! Illegal unit of measure (pt inserted).' Used as a
% <number>, which is what \setcounter uses it as, TeX coerces the dimension to
% scaled points, and 16383.99998pt is 1073741823sp exactly. That number is what
% is written here: it is what tex answers for this use, and the only use this
% layer has for \maxdimen is this one.
\def\maxdimen{1073741823}
%
% ---- plain.tex's allocators ----------------------------------------------
% plain.tex:
%
%     \def\alloc@#1#2#3#4#5{\global\advance\count1#1by\@ne
%       \ch@ck#1#4#2\allocationnumber=\count1#1
%       \global#3#5=\allocationnumber ...}
%     \def\newcount{\alloc@0\count\countdef\insc@unt}
%     \def\newdimen{\alloc@1\dimen\dimendef\insc@unt}
%     \def\newskip {\alloc@2\skip \skipdef \insc@unt}
%     \def\newtoks {\alloc@5\toks \toksdef \@cclvi}
%     \def\newbox  {\alloc@4\box  \chardef \insc@unt}
%
% The middle line is the one that cannot be carried: `\global#3#5=\allocationnumber'
% freezes the value of a count register into a NAME, and nothing freezes in this
% engine -- the four measurements in the counters section above say so at
% length, and this is the fifth thing they rule out.
%
% So the arithmetic is done before the run, in `src/latex/allocate.rs', which
% reads this file and the document's own for the twelve allocating commands and
% writes the \countdef, \dimendef, \skipdef, \toksdef and \chardef each name
% needs. By the time the mouth reaches the \newcount itself the name is already
% defined, so what is left for these to do is to CONSUME the name that follows
% -- which is what they do here. This is the same division as \newif and
% \newcounter: name arithmetic the expander cannot perform is performed ahead
% of the run and handed to it as definitions.
%
% A name allocate.rs could not resolve -- `\expandafter\newcount\csname c@#1\endcsname'
% -- is consumed here and stays undefined, which is loud at its first use rather
% than quietly sharing another allocation's register.
\def\newcount#1{}
\def\newdimen#1{}
\def\newskip#1{}
\def\newmuskip#1{}
\def\newtoks#1{}
\def\newbox#1{}
\def\newinsert#1{}
\def\newread#1{}
\def\newwrite#1{}
\def\newlanguage#1{}
\def\newfam#1{}
% latex.ltx (ltlength): \newcommand*\newlength[1]{\@ifdefinable#1{\newskip#1}}
% A LaTeX length is a SKIP, not a dimension. Written with a plain parameter so
% that both spellings a document uses -- \newlength\foo and \newlength{\foo} --
% consume exactly the name.
\def\newlength#1{}
%
% ---- plain.tex's category-code names -------------------------------------
% plain.tex:  \chardef\active=13
% A number with a name, and the one ifthen.sty stops on: `\catcode`\+\active'
% is `! Missing number, found \active.' without it. \chardef is what plain uses
% and what texrs has, so this is the definition rather than a stand-in for it.
\chardef\active=13
% latex.ltx:132.  \def\@makeother#1{\catcode`#1=12\relax}
\def\@makeother#1{\catcode`#1=12\relax}
% latex.ltx:1786-1787. Every character a file name or a \write argument might
% carry, made harmless. Carried whole.
\def\@sanitize{\@makeother\ \@makeother\\\@makeother\$\@makeother\&%
\@makeother\#\@makeother\^\@makeother\_\@makeother\%\@makeother\~}
%
% ---- what a package says when something is wrong -------------------------
% latex.ltx:8913 and its neighbours:
%
%     \gdef\PackageError#1#2#3{\GenericError{(#1)\@spaces...}{Package #1 Error: #2}...}
%
% SUBSTITUTION. \GenericError's body is \errmessage inside a group that makes
% `~' expand to a \typeout, and texrs has neither \errmessage nor \typeout --
% there is no error recovery for one to return from. What these are needed for
% is that they be DEFINED: calc.sty, and most packages, name them only inside a
% macro body that a correct document never reaches, and the load stopped on the
% name rather than on anything the document did. So each consumes exactly the
% arguments latex.ltx gives it and produces nothing.
%
% The cost is written down rather than hidden: a package that really does raise
% an error during its own load is now silent about it, where before it stopped
% the load and was reported by name. That is the wrong direction for a report
% whose whole point is to name what is missing, and it is why these are here
% rather than in `prelude.tex' -- the substitution is deliberate and this file
% is where the deliberate ones are recorded.
\def\GenericError#1#2#3#4{}
\def\GenericWarning#1#2{}
\def\GenericInfo#1#2{}
\def\PackageError#1#2#3{}
\def\PackageWarning#1#2{}
\def\PackageWarningNoLine#1#2{}
\def\PackageInfo#1#2{}
\def\PackageInfoNoLine#1#2{}
\def\ClassError#1#2#3{}
\def\ClassWarning#1#2{}
\def\ClassWarningNoLine#1#2{}
\def\ClassInfo#1#2{}
\def\@latex@error#1#2{}
\def\@latex@warning#1{}
\def\@latex@warning@no@line#1{}
\def\@latex@info#1{}
\def\@latex@info@no@line#1{}
\def\@font@warning#1{}
\def\@font@info#1{}
\def\typeout#1{}
\def\wlog#1{}
\def\MessageBreak{}
\def\@spaces{    }
%
% ---- \endinput -----------------------------------------------------------
% NOT DEFINED, on purpose, and the reason is worth recording because defining
% it as nothing was TRIED and cost thirteen books. \endinput is a PRIMITIVE: it
% ends the current file after the line it appears on. A macro that does nothing
% leaves the rest of the file to be read, and a .sty that reaches \endinput from
% inside a conditional -- to stop before code meant for a newer format -- then
% reads exactly the part it meant to skip. Measured: with `\def\endinput{}' the
% sweep fell from 228 documents to 215.
%
% What IS handled is the only shape where the two readings agree, and it is
% handled where a file is opened rather than here: `Lowerer::open_input' drops a
% trailing \endinput when nothing but comments follows it, which is what
% docstrip writes at the end of every generated .sty -- keyval.sty:87 and
% inputenc.sty:164 are both `\endinput' followed by the two `%% End of file'
% lines. There, ending the file and reading on to the end of it are the same
% thing. Everywhere else \endinput is still undefined, so a file that needs the
% real primitive is REPORTED rather than misread.
%
% ---- \@setfontsize -------------------------------------------------------
% latex.ltx:14295:
%
%     \def\@setfontsize#1#2#3{\@nomath#1%
%       \ifx\protect\@typeset@protect \let\@currsize#1\fi
%       \fontsize{#2}{#3}\selectfont}
%
% SUBSTITUTION, over the whole of NFSS. \fontsize and \selectfont are already
% stand-ins in `prelude.tex' -- BUGS.md records that Layout::size is one
% document-wide number, so every heading is set at body size and giving these a
% body would change nothing -- and \@nomath is a math-mode warning this engine
% has no \ifmmode to test. What is left of the definition is the three
% arguments, consumed.
%
% It is the last thing between article.cls and a load: the class's own final
% line reads size10.clo, and size10.clo's first definition is
% `\renewcommand\normalsize{\@setfontsize\normalsize\@xpt\@xiipt ...}'.
\def\@setfontsize#1#2#3{}
\def\@setsize#1#2#3#4{}
%
% ---- \pagenumbering -------------------------------------------------------
% latex.ltx (ltpage):
%
%     \def\pagenumbering#1{\global\c@page \@ne
%       \gdef\thepage{\csname @#1\endcsname\c@page}}
%
% Two substitutions, both forced by how a counter is reached here. \c@page is
% the register, and this layer reaches a counter by NAME through \cr@page (see
% the counters section above), so the reset is \setcounter. And `\@arabic\c@page'
% takes the register where `\arabic{page}' takes the name, so what \thepage is
% redefined to is the name-taking form -- \csname arabic\endcsname{page} builds
% exactly the \arabic this file defines, for whichever of the five styles the
% argument names.
%
% It is minimal.cls's only unmet name, and minimal.cls is the shortest real
% class there is: \ProvidesClass, two \setlength, \renewcommand\normalsize,
% this, and \pagestyle.
\def\pagenumbering#1{\setcounter{page}{1}%
  \@namedef{thepage}{\csname #1\endcsname{page}}}
%
% ---- appending to a macro -------------------------------------------------
% latex.ltx (ltdefns):
%
%     \long\def\g@addto@macro#1#2{\begingroup
%       \toks@\expandafter{#1#2}\xdef#1{\the\toks@}\endgroup}
%
% SUBSTITUTION, in the ONE step latex.ltx takes that this engine cannot: the
% \xdef. Nothing freezes into a macro body here -- the counters section above
% measures that four ways -- so `\xdef#1{\the\toks@}' writes the three tokens
% \the \toks@ rather than the list.
%
% What is written instead is the \expandafter form the same append has been
% written in since before \toks registers were cheap, and which etoolbox's
% \g@addto@macro-alikes still use:
%
%     \expandafter\def\expandafter#1\expandafter{#1#2}
%
% Three \expandafters, one per token between the \def and the macro that has to
% be expanded before the brace closes. Measured here: \def\hook{A} followed by
% two appends says ABC, where every shape built on \xdef says \the\toks@.
%
% The one difference from latex.ltx's, and it is worth stating because the name
% carries a `g': this append is LOCAL where latex.ltx's is global. Everything
% that appends here does it in a preamble, which is one group, so the two agree
% wherever this layer is actually used -- but an append inside a group is undone
% by the group, and latex.ltx's is not.
\long\def\g@addto@macro#1#2{\expandafter\def\expandafter#1\expandafter{#1#2}}
\long\def\@addto@macro#1#2{\expandafter\def\expandafter#1\expandafter{#1#2}}
% latex.ltx (ltfiles): \addto@hook appends to a TOKEN REGISTER rather than to a
% macro -- `\def\addto@hook#1#2{#1\expandafter{\the#1#2}}'. That form is refused
% here: `\toks3=\expandafter{\the\toks3 Y}' answers `! Missing { inserted.',
% because a toks assignment does not expand what follows the `='. amsmath opens
% with one. The macro append above is the nearest thing that works, and a hook
% this layer never runs is a hook whose contents cost nothing either way.
\def\addto@hook#1#2{}
%
% ---- names every file assumes ---------------------------------------------
% latex.ltx: \def\space{ }
\def\space{ }
% latex.ltx:16033 and its neighbours declare the list parameters with \newskip
% and \newdimen; `src/latex/allocate.rs' reads these lines and hands each one a
% register. \@listI is the top-level list's settings and \@listi is \let to it,
% which is the pair size10.clo's \normalsize restores and natbib reads.
\newskip\topsep
\newskip\partopsep
\newskip\itemsep
\newskip\parsep
\newskip\@topsep
\newskip\@topsepadd
\newdimen\leftmargin
\newdimen\rightmargin
\newdimen\listparindent
\newdimen\itemindent
\newdimen\labelwidth
\newdimen\labelsep
\newdimen\linewidth
\newdimen\leftmargini
\newdimen\leftmarginii
\newdimen\leftmarginiii
\newdimen\leftmarginiv
\newdimen\leftmarginv
\newdimen\leftmarginvi
\newcount\@listdepth
\newcount\@itempenalty
\newcount\@beginparpenalty
\newcount\@endparpenalty
\def\@listI{}
\let\@listi\@listI
\def\@listii{}
\def\@listiii{}
\def\@listiv{}
\def\@listv{}
\def\@listvi{}
% latex.ltx:9138, 10758-10759. In tex.web these three are token PARAMETERS of
% the engine (§230); latex.ltx declares them with \newtoks because LaTeX takes
% them over, and that is the declaration this layer can carry -- a token
% register is a store texrs has.
\newtoks\everypar
\newtoks\everymath
\newtoks\everydisplay
% latex.ltx:8790-8800. The scratch registers every package helps itself to.
\newcount\@tempcnta
\newcount\@tempcntb
\newdimen\@tempdima
\newdimen\@tempdimb
\newdimen\@tempdimc
\newskip\@tempskipa
\newskip\@tempskipb
\newtoks\@temptokena
\newbox\@tempboxa
\newif\if@tempswa
%
% ---- plain.tex's named registers and constants ----------------------------
% plain.tex allocates these by hand rather than through \newcount, and writes
% the numbers out:
%
%     \countdef\count@=255   \dimendef\dimen@=0   \skipdef\skip@=0
%     \toksdef\toks@=0       \chardef\@ne=1       \mathchardef\@m=1000
%     \newdimen\p@ \p@=1pt   \newdimen\z@ \z@=0pt
%
% Carried as written, register numbers and all, because the numbers ARE the
% definition -- plain reserves 0..9 in each file for exactly this, and
% `src/latex/allocate.rs' hands out from 10 upward so the two cannot collide.
\countdef\count@=255
\dimendef\dimen@=0
\dimendef\dimen@i=1
\dimendef\dimen@ii=2
\skipdef\skip@=0
\toksdef\toks@=0
\dimendef\p@=3
\p@=1pt
\dimendef\z@=4
\z@=0pt
\skipdef\z@skip=1
\z@skip=0pt plus0pt minus0pt
\chardef\@ne=1
\chardef\tw@=2
\chardef\thr@@=3
\chardef\sixt@@n=16
\chardef\@cclv=255
\mathchardef\@cclvi=256
\mathchardef\@m=1000
\mathchardef\@M=10000
\mathchardef\@MM=20000
% plain.tex: \newcount\m@ne \m@ne=-1.
%
% SUBSTITUTION, and a small one: a \chardef cannot hold a negative, and a
% \newcount here would be a register `src/latex/allocate.rs' hands out -- which
% would make THIS FILE depend on the allocation block being in front of it, and
% `texrs::latex::PRELUDE' is preloaded on its own by callers that want only the
% definitions. So the value is carried as a macro. Everything plain uses \m@ne
% for reads it as a number -- `\advance\count@\m@ne' -- and none of it assigns
% to it, which is the only thing a register would add.
\def\m@ne{-1}
% latex.ltx:1202-1203. The two keywords, as macros, which is what lets
% `\@plus2\p@' read as `plus 2pt' when the glue is really scanned.
\def\@minus{minus}
\def\@plus{plus}
%
% ---- expl3's syntax switch ------------------------------------------------
% l3names: \ExplSyntaxOn makes space and tab IGNORED, `~' a space, and `_' and
% `:' letters, so that expl3's own names read as control words; \ExplSyntaxOff
% puts all five back. Both are catcode changes and nothing else, which is
% exactly what this engine can do.
%
% The one line of the real pair not carried is `\endlinechar 32', which texrs
% has no \endlinechar parameter for. It matters only to code that relies on an
% end of line producing a space where a space is ignored, and the packages this
% reaches -- array.sty wraps one \def in the pair, xcolor.sty wraps two lines --
% do not.
%
% It is NOT expl3. A file that uses expl3's own functions still stops on the
% first `\tl_new:N', by name, as it should.
%
% l3names writes the pair as bare CHARACTER CODES -- `\catcode 9 = 9 \scan_stop:'
% and so on -- and that is how they are carried here, because ``\^^I' is not a
% number this engine's scanner reads: `\catcode`\^^I=9' answers `! Missing
% number, treated as zero.' and `\catcode9=9' works. Nine is the tab either way.
\def\ExplSyntaxOn{\catcode9=9\relax\catcode32=9\relax\catcode34=12\relax
  \catcode58=11\relax\catcode94=7\relax\catcode95=11\relax
  \catcode124=12\relax\catcode126=10\relax}
\def\ExplSyntaxOff{\catcode9=10\relax\catcode32=10\relax\catcode34=12\relax
  \catcode58=12\relax\catcode94=7\relax\catcode95=8\relax
  \catcode124=12\relax\catcode126=13\relax}
%
% ---- NFSS, which this engine has no fonts for ----------------------------
% latex.ltx (ltfssbas, ltfssdcl, ltoutenc). Every one of these DECLARES a font,
% an encoding or a math symbol, and the whole of what they declare is read by a
% stomach choosing a face -- which `src/typeset.rs' does from the document's own
% \setmainfont rather than from NFSS's tables. So each consumes the arguments
% latex.ltx gives it and produces nothing, and the note BUGS.md already carries
% about one type size covers what that costs.
%
% They are here rather than in `prelude.tex' because the arities are
% latex.ltx's and getting one wrong is worse than leaving the name undefined:
% a stand-in that eats one argument too few leaves the rest in the text.
\def\DeclareFontFamily#1#2#3{}
\def\DeclareFontShape#1#2#3#4#5#6{}
\def\DeclareFontEncoding#1#2#3{}
\def\DeclareFontSubstitution#1#2#3#4{}
\def\DeclareFontEncodingDefaults#1#2{}
\def\DeclareErrorFont#1#2#3#4#5{}
\def\DeclarePreloadSizes#1#2#3#4#5{}
\def\DeclareSymbolFont#1#2#3#4#5{}
\def\SetSymbolFont#1#2#3#4#5#6{}
\def\DeclareSymbolFontAlphabet#1#2{}
\def\DeclareMathAlphabet#1#2#3#4#5{}
\def\SetMathAlphabet#1#2#3#4#5#6{}
\def\DeclareMathVersion#1{}
\def\DeclareMathSymbol#1#2#3#4{}
\def\DeclareMathDelimiter#1#2#3#4#5#6{}
\def\DeclareMathRadical#1#2#3#4#5{}
\def\DeclareMathAccent#1#2#3#4{}
\def\DeclareTextCommand#1#2{\@ifnextchar[\@dtc@opt\@dtc@plain}
\def\@dtc@opt[#1]{\@ifnextchar[\@dtc@opt@dflt\@dtc@opt@body}
\def\@dtc@opt@dflt[#1]#2{}
\def\@dtc@opt@body#1{}
\def\@dtc@plain#1{}
\let\ProvideTextCommand\DeclareTextCommand
\let\DeclareTextCommandDefault\@gobbletwo
\let\ProvideTextCommandDefault\@gobbletwo
\def\DeclareTextSymbol#1#2#3{}
\def\DeclareTextSymbolDefault#1#2{}
\def\DeclareTextAccent#1#2#3{}
\def\DeclareTextAccentDefault#1#2{}
\def\DeclareTextComposite#1#2#3#4{}
\def\DeclareTextCompositeCommand#1#2#3#4{}
\def\DeclareTextFontCommand#1#2{}
\def\usefont#1#2#3#4{}
\def\fontencoding#1{}
\def\fontfamily#1{}
\def\fontseries#1{}
\def\fontshape#1{}
\def\linespread#1{}
\def\DeclareOldFontCommand#1#2#3{}
%
% ---- file writing --------------------------------------------------------
% tex.web §1350: \write is a whatsit the SHIPPER expands, and \immediate makes
% it happen now instead. There is no shipper here and no \openout, so both are
% consumed. \write's argument is a stream number and a group, which is what the
% two parameters take; the stream number is always a name `\newwrite' made,
% so it arrives as one token.
\def\immediate{}
\def\write#1#2{}
\def\openout#1{}
\def\closeout#1{}
\def\@gobbletwo#1#2{}
